%% kapa: exam=1 round=33 year=2022 period=1 form=A subject=real_estate subject_ko=부동산학원론 seq=24 number=104
%% source: Q-Net 감정평가사 기출문제 (cst00309) artlSeq=5209997 file=raw/1008/5209997/2022년도 제33회 감정평가사 1차시험 1교시(민법,경제학,부동산학).pdf
%% conversion: kordoc 4.24.51 -> tools/gen1.py (auto); review=proofread-sonnet; flags=none
\begin{question}{104}
건물의 관리방식에 관한 설명으로 옳은 것은?
\begin{choices}
  \item 위탁관리방식은 부동산관리 전문업체에 위탁해 관리하는 방식으로 대형건물의 관리에 유용하다.
  \item 혼합관리방식은 필요한 부분만 일부 위탁하는 방식으로 관리자들간의 협조가 긴밀하게 이루어진다.
  \item 자기관리방식은 관리업무의 타성(惰性)을 방지할 수 있다.
  \item 위탁관리방식은 외부 전문가가 관리하므로 기밀 및 보안 유지에 유리하다.
  \item 혼합관리방식은 관리문제 발생시 책임소재가 명확하다.
\end{choices}
\end{question}
\answer{1}
